\chapter{Saint-Sa\"{e}ns, Camille} % (fold)
\label{cha:saintsaens}

\href{https://en.wikipedia.org/wiki/List_of_compositions_by_Camille_Saint-Saëns}{List of works} on Wikipedia.





\section{Orchestral}

\noindent{Comic Opera Overture in E minor (c.~1850)} \\
\noindent{Scherzo in A major (c.~1850)} \textit{Orchestre Philharmonique Royal de Li\`{e}ge, Jean-Jacques Kantorow}\marginnote{max} \\
\noindent{Symphony in A major (c.~1850), R.~159} \\
\noindent{Symphony No. 1 in E$\flat$ major (1853), Op.~2} R.~161 \textit{Orchestre Philharmonique Royal de Li\`{e}ge, Jean-Jacques Kantorow}\marginnote{09/2021, max} \\
\noindent{Overture to an Unfinished Comic Opera in G major (1854), R.~162} \\
\noindent{Symphony in F major "Urbs Roma" (1856), R.~163} \\
\noindent{Tarantelle in A minor for flute, clarinet, and orchestra (1857), Op.~6, R.~183} \\
\noindent{Violin Concerto No. 2 in C major (1858), Op.~58, R.~184} \\
\noindent{Piano Concerto No. 1 in D major (1858), Op.~17, R.~185} \textit{Gabriel Tacchino, Orchestra of Radio Luxembourg, Louis de Froment}\marginnote{high} \\
\noindent{Violin Concerto No. 1 in A major (1859), Op.~20, R.~186} \\
\noindent{Symphony No. 2 in A minor (1859), Op.~55} \textit{Orchestre Philharmonique Royal de Li\`{e}ge, Jean-Jacques Kantorow}\marginnote{09/2021, max} \\
\noindent{Suite for cello and orchestra (1862, 1919), Op.~16b, R.~211} \\
\noindent{Introduction and Rondo Capriccioso in A minor for violin and orchestra (1863), Op.~28, R.~188} \\
\noindent{Suite in D major (1863), Op.~49, R.~165} \\
\noindent{Spartacus, Overture in E$\flat$ major (1863), R.~166} \\
\noindent{Rhapsodie bretonne (1866, 1891), Op.~7b, R.~177} \\
\noindent{Piano Concerto No. 2 in G minor (1868), Op.~22, R.~190} \textit{Gabriel Tacchino, Orchestra of Radio Luxembourg, Louis de Froment}\marginnote{high} \\
\noindent{Piano Concerto No. 3 in E$\flat$ major (1869), Op.~29, R.~191} \textit{Gabriel Tacchino, Orchestra of Radio Luxembourg, Louis de Froment}\marginnote{high} \\
\noindent{Heroic March in E$\flat$ major (1871), Op.~34, R.~168} \\
\noindent{Romance in D$\flat$ major for flute and piano or orchestra (1871), Op.~37, R.~192} \\
\noindent{Evening Song (1872)} \\
\noindent{The Spinning Wheel of Omphale in A major (1872), Op.~31, R.~169} \textit{Les Si\`{e}cles, Fran\c{c}ois-Xavier Roth}\marginnote{12/2023, max} \\
\noindent{Cello Concerto No. 1 in A minor (1872), Op.~33, R.~193} \textit{Bruno Philippe, Frankfurt RSO, Christoph Eschenbach}\marginnote{01/2001, quality} \\
\noindent{Phaéton in C major (1873), Op.~39, R.~170} \textit{Les Si\`{e}cles, Fran\c{c}ois-Xavier Roth}\marginnote{12/2023, max} \\
\noindent{Romance in F major for horn (or cello) and orchestra (1874), Op.~36, R.~195} \\
\noindent{Danse macabre, Symphonic Poem in G minor (1874), Op.~40, R.~171} \textit{Les Si\`{e}cles, Fran\c{c}ois-Xavier Roth}\marginnote{12/2023, max} \\
\noindent{Romance in C major for violin and orchestra (1874), Op.~48, R.~196} \\
\noindent{Piano Concerto No. 4 in C minor (1875), Op.~44, R.~197} \textit{Daniel Roth, Les Si\`{e}cles, Fran\c{c}ois-Xavier Roth}\marginnote{max} \\
\noindent{La jeunesse d'Hercule (1877), Op.~50, R.~172} \textit{Les Si\`{e}cles, Fran\c{c}ois-Xavier Roth}\marginnote{12/2023, max} \\
\noindent{Cadenzas for Beethoven's Piano Concerto No. 4 (c. 1878)} \\
\noindent{Suite algérienne (1880), Op.~60, R.~173} \\
\noindent{Violin Concerto No. 3 in B minor (1880), Op.~61, R.~198} \textit{Rudolf Koelman, Sinfonietta Schaffhausen, Paul K. Haug}\marginnote{max}  \\
\noindent{Concert Piece in E minor for violin and orchestra (1880), Op.~62} \\
\noindent{Une nuit à Lisbonne (1880), Op.~63, R.~175} \\
\noindent{La jota aragonese (1880), Op.~64, R.~174} \\
\noindent{Hymn to Victor Hugo (1881), Op.~69} \\
\noindent{Allegro appassionato in C$\sharp$ minor for piano and orchestra (1884), Op.~70, R.~37} \\
\noindent{Rhapsody of the Auvergne in C major for piano and orchestra (1884), Op.~73} \\
\noindent{Romance in E major for horn (or cello) and orchestra (1885), Op.~67, R.~189} \\
\noindent{Wedding Cake, Caprice-valse in A$\flat$ major for piano and orchestra (1885), Op.~76, R.~124} \\
\noindent{Symphony No. 3 "Organ Symphony" in C minor (1886), Op.~78, R.~176} \textit{Daniel Roth, Les Si\`{e}cles, Fran\c{c}ois-Xavier Roth}\marginnote{max} \\
\noindent{Havanaise in E major for violin and orchestra (1887), Op.~83, R.~202} \\
\noindent{Concert Piece in F minor for horn and orchestra (1887), Op.~94} \\
\noindent{Africa in G minor for piano and orchestra (1891), Op.~89, R.~204} \textit{Gabriel Tacchino, Orchestra of Radio Luxembourg, Louis de Froment}\marginnote{high} \\
\noindent{Sarabande and Rigaudon in E major (1892), Op.~93} \\
\noindent{Piano Concerto No. 5 "Egyptian" in F major (1896), Op.~103, R.~205} \textit{Gabriel Tacchino, Orchestra of Radio Luxembourg, Louis de Froment}\marginnote{high} \\
\noindent{Cadenzas for Beethoven's Violin Concerto (1900)} \\
\noindent{Coronation March in E$\flat$ major (1902), Op.~117, R.~180} \\
\noindent{Cello Concerto No. 2 in D minor (1902), Op.~119} \\
\noindent{Caprice andalou in G major for violin and orchestra (1904), Op.~122, R.~207} \\
\noindent{Andante extrait d'un concerto pour piano for violin and orchestra (or piano) (1906)} \\
\noindent{3 Symphonic Pictures after "La Foi" (1908), Op.~130} \\
\noindent{La muse et le poète in E minor for violin, cello, and orchestra (or piano) (1910), Op.~132, R.~208} \\
\noindent{Ouverture de fête (1910), Op.~133, R.~181} \\
\noindent{Cadenzas for Mozart's Piano Concerto No. 22, K.482 (1911), R.~281} \\
\noindent{Cadenzas for Mozart's Piano Concerto No. 24, K.491 (1911), R.~279} \\
\noindent{Hail! California (1915)} \\
\noindent{Concert Piece in G major for harp and orchestra (1918), Op.~154, R.} \\
\noindent{Cyprès et lauriers in D minor for organ and orchestra (1919), Op.~156, R.~210} \\
\noindent{Odelette in D major for flute and orchestra (1920), Op.~162, R.~212} \\

\section{Chamber}

\noindent{Le carnaval des animaux, grande fantaisie zoologique (1886), R.~125} \textit{Les Si\`{e}cles, Fran\c{c}ois-Xavier Roth}\marginnote{12/2023, max}\\ 


\section{Film}

\noindent{L'assassinat du duc de Guise (1908), Op.~128, R.~331} \textit{Les Si\`{e}cles, Fran\c{c}ois-Xavier Roth}\marginnote{12/2023, max}\\